\section{Introduction}\label{sec:intro}

Move is a bytecode language with first-class references over local
and globally-stored values. Each function's bytecode passes through
a static reference-safety verifier, which approves the function only
if, on every reachable execution path, references stay in scope, do
not violate type safety through reads or writes, do not overlap as
mutable arguments at call boundaries, and do not allow an immutable
view to observe a value rewritten through another handle. The
verifier rejects bytecode that violates these rules.

We are interested in enforcing the same rules \emph{at execution
time}. There are several reasons one would want this:
\begin{itemize}
\item Fuzzing and mutated-bytecode harnesses deliberately exercise
inputs the verifier would reject, to find bugs in everything
downstream.
\item Dynamic dispatch and native code escape static analysis: the
verifier cannot see what happens through a dispatched function
pointer or inside a Rust-implemented native, but the reference
rules still need to hold across those boundaries.
\item The static check has to be sound, so it sometimes rejects
programs that would not actually violate any rule on the executions
they take. A runtime check sees the executions and can be more
precise.
\end{itemize}

The following Move fragment illustrates the kind of subtle case at
stake:
\begin{lstlisting}
enum E { V1(T1), V2(T2) }
let v: E = ...;
let inner: &mut T1 = match &mut v { V1(p) => p };
... // some path mutates v to be V2
*inner    // unsafe
\end{lstlisting}
The reference \texttt{inner} is in scope. Its lifetime has not ended.
Its declared type is unchanged. But the program has switched
\texttt{v} from \texttt{V1} to \texttt{V2}, so the substructure that
\texttt{inner} pointed into is no longer there. Dereferencing
\texttt{inner} now reads memory that has been recycled into a value
of an unrelated type.

A sound static analysis catches this only conservatively: it does
not know which branch will run, so it must assume the worst. By the
time \texttt{*inner} executes, the runtime has already observed which
branch ran, which mutation occurred, and which reference was derived
from which. With the right slice of that information recorded, we
can answer ``is this reference still safe to use?'' precisely. The
challenge is to find a representation that is small enough to be
cheap to maintain, and uniform enough to cover the various events
that can affect a reference.

\paragraph{This paper.}
We organize live references by the path of borrows that produced
them, using a hierarchical structure we call an \emph{Access Path
Tree} (\apt). Each live reference is attached to a node of some APT,
with edges labelled by field projections, variant projections, and
a single label for vector indices. A single rule decides when a
reference becomes unusable: when an event destructively changes the
substructure reachable through some node, every reference attached
to a related node is marked \emph{poisoned} for the rest of the
execution. Destructive writes, moves and overwrites, freezes,
variant mutations, calls, returns, and opaque effects all fit this
pattern; what varies between them is which set of related nodes is
affected and which mutability filter applies. Calls additionally
lock subtrees so that two mutable arguments cannot overlap.

\paragraph{Contributions.}
\begin{itemize}
\item A runtime discipline for Move's reference-safety rules,
  organized around access path trees and one invalidation rule
  (Sects.~\ref{sec:apt}, \ref{sec:semantics}).
\item A small-step presentation in which the various
  reference-affecting events are instances of the same rule
  (Sect.~\ref{sec:semantics}).
\item A discussion of what the runtime trace buys over a join-based
  static check (Sect.~\ref{sec:precision}) and how the discipline
  relates to per-location dynamic alias models such as Stacked
  Borrows and Tree Borrows~\cite{stackedborrows,treeborrows}
  (Sect.~\ref{sec:related}).
\end{itemize}
